% Work/MAT – učitelská příprava pro 1. ročník; přepracované pořadí bloku MNZ.
% Hodina 20: Opakování množin a intervalů (přesun z původní 18). 45 minut povinného jádra;
% všechno za oddílem "Volitelná rezerva" je nepovinné rozšíření / domácí práce.
% Předchozí hodiny: 12 Množiny; 13 Podmnožiny a operace; 14 Vennovy diagramy;
% 15 Vlastnosti a volitelně mohutnost; 16 Operace s intervaly;
% 17 – dvě množiny; 18 – tři množiny; 19 – procvičení slovních úloh.
% Vlastní zdroje: ucebnice/01_Zakladni_poznatky/Mnoziny/Množiny.tex;
% ucebnice/01_Zakladni_poznatky/Ciselne_obory/Číselné obory.tex
% (oddíl Intervaly); původní příprava 20_MNZ_OpakovaniMnozin.tex
% a přepracované přípravy 16 až 19.
% Externí metodika: Realisticky.cz 010309 Množiny – shrnutí (doložený
% veřejně indexovaný text: shrnout pojmy, typické chyby a užitečné postupy),
% 010308 Intervaly; 010304 Vennovy diagramy; 010306 a 010307 slovní úlohy.
% Konkrétní příklady a jejich řešení jsou vlastní tvorba, nikoli doslovný přepis.
% Kontrola: úplnost výsledků, krajní body, součty v diagramech, průběh 45 minut.
\documentclass[11pt,a4paper]{article}
% Spolecny styl pro ucitelske pripravy MAT (Work/MAT)
% Verze 1.1, 2026-10-23
\NeedsTeXFormat{LaTeX2e}
\ProvidesPackage{priprava}[2026/10/23 v1.1 Sablona ucitelskych priprav MAT]

% Zakladni typografie a cestina
\RequirePackage[T1]{fontenc}
\RequirePackage[utf8]{inputenc}
\RequirePackage[czech]{babel}
\RequirePackage{lmodern}
\RequirePackage{microtype}
\RequirePackage{amsmath,amssymb}
\RequirePackage{xcolor}
\RequirePackage{enumitem}
\RequirePackage[most]{tcolorbox}
\RequirePackage{tabularx,booktabs,array}
\RequirePackage{fancyhdr}
\RequirePackage{needspace}
\RequirePackage{geometry}
\RequirePackage[hidelinks]{hyperref}

\geometry{a4paper,top=17mm,bottom=18mm,left=23mm,right=23mm}
\definecolor{notegray}{gray}{0.39}
\setlength{\parindent}{0pt}
\setlength{\parskip}{2.5pt plus 1pt minus 0.5pt}
\setlist[itemize]{topsep=2pt,itemsep=1pt,parsep=0pt,leftmargin=1.45em}
\setlist[enumerate]{topsep=2pt,itemsep=2pt,parsep=0pt,leftmargin=1.65em}
\setcounter{secnumdepth}{0}
\renewcommand{\arraystretch}{1.12}

% Zahlavi: vlevo rocnik | globalni tema, vpravo cislo: nazev hodiny
% Zapati zustava zcela prazdne.
\pagestyle{fancy}
\fancyhf{}
\renewcommand{\headrulewidth}{0.35pt}
\renewcommand{\footrulewidth}{0pt}
\setlength{\headheight}{14pt}
\newcommand{\nastavPripravu}[4]{%
  \fancyhead[L]{\small #1\enspace\textbar\enspace #2}%
  \fancyhead[R]{\small #3: #4}%
}

% Minimalisticke cernobile boxy. Kazda definice zvlast.
\tcbset{commonlesson/.style={enhanced,breakable=false,colback=white,colframe=black,
  colbacktitle=white,coltitle=black,boxrule=0.45pt,
  arc=0.8mm,outer arc=0.8mm,left=3mm,right=3mm,top=1.6mm,bottom=1.6mm,
  toptitle=1mm,bottomtitle=0.7mm,fonttitle=\bfseries,
  attach boxed title to top left={xshift=2.5mm,yshift=-1.15mm},
  boxed title style={colback=white,colframe=white,boxrule=0pt,left=1mm,right=1mm}}}
\newtcolorbox{zadaniBox}[1][]{commonlesson,title={Zadání},#1}
\newtcolorbox{prikladBox}[1][]{commonlesson,title={Příklad},#1}
\newtcolorbox{definiceBox}[1][]{commonlesson,title={Definice},#1}

% Spolecny box pro zadani a jeho reseni, bez mezery mezi castmi.
% Pouziti:
%   \begin{prikladReseniBox}
%     Zadani ulohy.
%     \tcblower
%     \begin{align*} ... \end{align*}
%   \end{prikladReseniBox}
% Spodni cast ma jemny sedy podklad a delici linku; lze v ni sazet
% vice radku matematiky. Obe casti maji jen jeden spolecny ramecek.
\tcbset{solutionlesson/.style={
  commonlesson,
  enhanced jigsaw,
  breakable,
  colbacklower=black!3,
  segmentation style={solid,draw=black!25,line width=0.35pt},
  fontlower=\small\color{notegray},
  before lower={\noindent\textbf{Řešení:}\enspace},
  leftlower=3mm,rightlower=3mm
}}
\newtcolorbox{prikladReseniBox}[1][]{solutionlesson,title={Příklad},#1}
\newtcolorbox{zadaniReseniBox}[1][]{solutionlesson,title={Zadání},#1}

% Pomucky pro autora pripravy: sedive poznamky; \reseni pro starsi pripravy
\newcommand{\poznamka}[1]{\par\begingroup\small\color{notegray}\itshape
  \textnormal{\textbf{Poznámka pro učitele.}}\enspace #1\par\endgroup}
\newcommand{\reseni}[1]{\par\begingroup\small\color{notegray}
  \textbf{Očekávané řešení:}\enspace #1\par\endgroup}
\newcommand{\blok}[1]{\needspace{5\baselineskip}\subsection*{#1}}
\newcommand{\vyrok}[1]{\uv{#1}}

\usepackage{tikz}
\usetikzlibrary{calc}
\nastavPripravu{1. ročník}{Množiny}{20}{Opakování množin a intervalů}
\begin{document}

\blok{Dnes žádný nový pojem: umět poznat a opravit chybu}
Navazujeme na množiny, Vennovy diagramy, množinové operace a intervaly.
Cílem není znovu opsat definice, ale \textbf{samostatně zvolit postup,
zdůvodnit výsledek a zkontrolovat jej}. Při řešení pracujeme v uvedené
základní množině $U$; pokud není řečeno jinak, intervaly chápeme v $\mathbb R$.

\poznamka{Časový rámec pro 45 minut:
D0 (4 min), S1 (8 min), S2 (8 min), S3 (11 min), S4 (11 min),
závěr (3 min). Součástí časů je společná kontrola. Při pomalejším tempu
ponech S3 c) a S4 poslední otázku na domácí dokončení.
Spodní části boxů s řešením odkrývej až po pokusu žáků.}

\begin{zadaniReseniBox}[title={D0 -- Čtyři chyby na rozjezd},breakable=false]
Rozhodni, zda jsou tvrzení správná. Chybné zápisy oprav a krátce vysvětli.
\begin{enumerate}[label=\alph*)]
\item $\{0\}=\varnothing$.
\item $\{1,2\}\subseteq\{2,3\}$.
\item $(-2;1)\cup(1;4)=(-2;4)$.
\item $A\setminus B=B\setminus A$ pro libovolné množiny $A,B$.
\end{enumerate}
\tcblower
\begin{enumerate}[label=\alph*)]
\item Ne. $\{0\}$ obsahuje jeden prvek, $\varnothing$ žádný.
\item Ne. Číslo $1$ je v první množině, ale ne ve druhé.
\item Ne. Bod $1$ není v žádném intervalu, proto jej nelze do sjednocení přidat.
\item Ne. Např. pro $A=\{1,2\}$, $B=\{2\}$ je
$A\setminus B=\{1\}$, ale $B\setminus A=\varnothing$.
\end{enumerate}
\end{zadaniReseniBox}

\blok{1. Když se setkají výčet, vlastnost a podmnožina}
\begin{zadaniReseniBox}[title={S1 -- Jedna úloha, několik základních pojmů}]
Nechť $U=\{-3,-2,-1,0,1,2,3,4,5\}$,
$A=\{x\in\mathbb Z\mid-2\le x<3\}$ a $B=\{-1,1,3\}$.
\begin{enumerate}[label=\alph*)]
\item Vypiš $A$. Je $B\subseteq A$? Své rozhodnutí zdůvodni.
\item Urči $A\cap B$, $A\cup B$, $A\setminus B$ a $B\setminus A$.
\item Urči doplněk $(A\cup B)'_U$ a počet prvků $|A\cup B|$.
\end{enumerate}
\tcblower
\begin{enumerate}[label=\alph*)]
\item $A=\{-2,-1,0,1,2\}$; $B\nsubseteq A$, protože $3\in B$,
ale $3\notin A$.
\item $A\cap B=\{-1,1\}$; $A\cup B=\{-2,-1,0,1,2,3\}$; \\
$A\setminus B=\{-2,0,2\}$; $B\setminus A=\{3\}$.
\item $(A\cup B)'_U=\{-3,4,5\}$ a $|A\cup B|=6$.
Také $|A|+|B|-|A\cap B|=5+3-2=6$.
\end{enumerate}
\end{zadaniReseniBox}
\poznamka{Diagnostika: odlišuj zápis $B\nsubseteq A$ od $3\notin A$.
Pro doplněk vždy explicitně požaduj otázku \uv{Vzhledem k jakému $U$?}
Svislé čáry v $|A|$ označují mohutnost, ne absolutní hodnotu.}

\blok{Pět kontrolních otázek, které předcházejí chybám}
\begin{center}
\begin{tabularx}{\linewidth}{@{}lX@{}}
\toprule
Téma & Co zkontroluji ještě před odevzdáním?\\ \midrule
Zadání množiny & Zvolil jsem správný číselný obor a krajní hodnoty?\\
Podmnožina & Prověřil jsem \emph{každý} prvek, nebo mám protipříklad?\\
Operace & Čtu správně \uv{a současně}, \uv{nebo}, \uv{ale ne}?\\
Interval & Patří oba krajní body do výsledku? Nechybí jediný bod?\\
Slovní úloha & Dávají součty nepřekrývajících se oblastí celý soubor?\\
\bottomrule
\end{tabularx}
\end{center}
\poznamka{Přehled je pomůcka k samostatné kontrole, nikoli další látka k opisování.}

\blok{2. Z Vennova diagramu přečti každou oblast}
\begin{zadaniReseniBox}[title={S2 -- Rozdělení prvků a složená operace},breakable=false]
V $U=\{1,2,\ldots,10\}$ jsou $A$ sudá čísla a $B$ násobky tří.
\begin{enumerate}[label=\alph*)]
\item Zakresli $A$, $B$ do Vennova diagramu a vyplň všechny čtyři oblasti.
\item Urči $(A\cup B)'_U$, $(A\setminus B)\cup(B\setminus A)$
a počet prvků v $A\cap B$.
\end{enumerate}
\tcblower
$A=\{2,4,6,8,10\}$, $B=\{3,6,9\}$.
Jen $A$: $\{2,4,8,10\}$; obě: $\{6\}$; jen $B$: $\{3,9\}$;
mimo obě: $\{1,5,7\}$.
\begin{center}
\begin{tikzpicture}[scale=.78]
 \draw[thick] (-3.2,-1.72) rectangle (3.2,1.86);
 \draw[thick] (-.82,0) circle(1.31);
 \draw[thick] (.82,0) circle(1.31);
 \node at (-1.43,1.52) {$A$}; \node at (1.43,1.52) {$B$};
 \node[anchor=north east] at (3.03,1.7) {$U$};
 \node[align=center] at (-1.30,0) {$2,4$\\$8,10$};
 \node at (0,0) {$6$};
 \node at (1.18,0) {$3,9$};
 \node at (2.25,-1.26) {$1,5,7$};
\end{tikzpicture}
\end{center}
$(A\cup B)'_U=\{1,5,7\}$;
$(A\setminus B)\cup(B\setminus A)=\{2,3,4,8,9,10\}$;
$|A\cap B|=1$.
\textbf{Kontrola:} $4+1+2+3=10$; žádný prvek není dvakrát.
\end{zadaniReseniBox}
\poznamka{Sleduj, zda žáci nezapíšou společnou šestku do obou kruhů dvakrát.
Požádej o slovní popis výrazu $(A\setminus B)\cup(B\setminus A)$:
\uv{právě v jedné z obou množin}.}

\blok{3. Intervaly: rozhodující jsou krajní body}
\begin{prikladReseniBox}[title={S3a -- Překryv a odebrání části intervalu}]
Pro $I=\langle-3;2)$ a $J=(0;5\rangle$ nakresli společnou
číselnou osu a urči $I\cap J$, $I\cup J$, $I\setminus J$ a $J\setminus I$.
\tcblower
\begin{center}
\begin{tikzpicture}[x=.84cm,y=.64cm]
 \draw[->,thin] (-4.3,0)--(6.1,0) node[right] {$x$};
 \foreach \x in {-3,0,2,5}{\draw (\x,.12)--(\x,-.12) node[below=3pt] {$\x$};}
 \node[left] at (-4.25,1) {$I$};
 \draw[line width=1.6pt] (-3,1)--(2,1);
 \fill (-3,1) circle(2.65pt); \draw[fill=white] (2,1) circle(2.65pt);
 \node[left] at (-4.25,2) {$J$};
 \draw[line width=1.6pt] (0,2)--(5,2);
 \draw[fill=white] (0,2) circle(2.65pt); \fill (5,2) circle(2.65pt);
 \node[left] at (-4.25,3) {$I\cap J$};
 \draw[line width=1.6pt] (0,3)--(2,3);
 \draw[fill=white] (0,3) circle(2.65pt); \draw[fill=white] (2,3) circle(2.65pt);
\end{tikzpicture}
\end{center}
\[
\begin{aligned}
I\cap J&=(0;2),& I\cup J&=\langle-3;5\rangle,\\
I\setminus J&=\langle-3;0\rangle,&J\setminus I&=\langle2;5\rangle.
\end{aligned}
\]
\textbf{Zkouška:} nula patří do $I$, nikoli do $J$, takže zůstává v $I\setminus J$.
Dvojka patří do $J$, nikoli do $I$, proto je v $J\setminus I$.
\end{prikladReseniBox}

\begin{zadaniReseniBox}[title={S3b--c -- Dvě záludnosti bez kreslení hotového řešení}]
\begin{enumerate}[label=\alph*),start=2]
\item $K=(-2;1\rangle$, $L=(1;4)$. Urči $K\cap L$ a $K\cup L$.
Vysvětli, proč může sjednocení tvořit jeden interval i při prázdném průniku.
\item $M=(-2;1)$, $N=(1;4)$. Urči $M\cap N$, $M\cup N$.
Které jediné číslo brání sloučení do jednoho intervalu?
\end{enumerate}
\tcblower
\begin{enumerate}[label=\alph*),start=2]
\item $K\cap L=\varnothing$ a $K\cup L=(-2;4)$.
Bod $1$ leží v $K$; čísla větší než $1$ leží v $L$.
Mezi $-2$ a $4$ proto ve sjednocení nic nechybí.
\item $M\cap N=\varnothing$, ale $M\cup N=(-2;1)\cup(1;4)$.
Číslo $1$ není v žádné množině; nelze je do výsledku doplnit.
\end{enumerate}
\end{zadaniReseniBox}

\blok{4. Slovní úloha: od textu k počtům}
\begin{prikladReseniBox}[title={S4 -- Dva školní kroužky a neznámý průnik},breakable=false]
Ve skupině $30$ studentů navštěvuje matematický seminář $17$,
programátorský kroužek $15$ a \emph{ani jeden} z nich $6$ studentů.
Urči: a) počet v obou skupinách; b) pouze v první; c) pouze ve druhé;
d) počet navštěvujících \emph{právě jeden} a \emph{nejvýše jeden} kroužek.
Sestav Vennův diagram s počty a ověř všechny údaje.
\tcblower
Označme $A$ seminář, $B$ programování.
Alespoň jeden navštěvuje $30-6=24$ studentů.
\[
|A\cap B|=|A|+|B|-|A\cup B|=17+15-24=8.
\]
Pouze seminář $17-8=9$, pouze programování $15-8=7$.
Právě jeden: $9+7=16$; nejvýše jeden: $9+7+6=22$.
\begin{center}
\begin{tikzpicture}[scale=.79]
 \draw[thick] (-2.7,-1.52) rectangle (2.9,1.63);
 \draw[thick] (-.8,0) circle(1.2); \draw[thick] (.8,0) circle(1.2);
 \node at (-1.35,1.4) {$A$}; \node at (1.35,1.4) {$B$};
 \node[anchor=north east] at (2.75,1.5) {$U$};
 \node at (-1.1,0) {\Large $9$}; \node at (0,0) {\Large $8$};
 \node at (1.12,0) {\Large $7$}; \node at (2.32,-1.13) {\Large $6$};
\end{tikzpicture}
\end{center}
Kontrola: $9+8+7+6=30$, $9+8=17$, $8+7=15$.
\end{prikladReseniBox}

\textbf{Výstupní lístek (3 min):} Napiš jednu věc, kterou vždy
kontroluješ u \emph{intervalů}, a jednu u \emph{slovních úloh na množiny}.
Přidej vlastní krátký příklad rozdílu množin.

\poznamka{Očekávej např. kontrolu zahrnutí krajních bodů a kontrolu
součtu disjunktních oblastí s celkovým počtem. Výstupní lístek ukáže,
co je třeba ještě procvičit. Tím končí povinné 45minutové jádro.}

\blok{Volitelná rezerva: tři množiny a výraz se závorkami}
Tato část \textbf{nepatří do 45minutového jádra}. Můžeš ji vybrat pro rychlejší
žáky, dobrovolný domácí úkol nebo další procvičovací hodinu.

\begin{prikladReseniBox}[title={R1 -- Osm oblastí tří množin}]
Ve skupině $40$ lidí platí $|A|=22$, $|B|=18$, $|C|=16$,
$|A\cap B|=8$, $|A\cap C|=7$, $|B\cap C|=6$ a
$|A\cap B\cap C|=3$. Vypočti všech osm oblastí diagramu,
počet v právě dvou skupinách a počet mimo všechny.
\tcblower
Nejprve společný střed: $3$.
Přesně dvě skupiny: $AB$ bez $C$: $8-3=5$;
$AC$ bez $B$: $7-3=4$; $BC$ bez $A$: $6-3=3$.
Pouze $A$: $22-5-4-3=10$;
pouze $B$: $18-5-3-3=7$;
pouze $C$: $16-4-3-3=6$.
\begin{center}
\begin{tikzpicture}[scale=.84]
 \draw[thick] (-2.95,-2.57) rectangle (3.02,2.05);
 \draw[thick] (-.82,.37) circle(1.4);
 \draw[thick] (.82,.37) circle(1.4);
 \draw[thick] (0,-.79) circle(1.4);
 \node at (-1.72,1.67) {$A$}; \node at (1.72,1.67) {$B$};
 \node at (0,-2.39) {$C$}; \node[anchor=north east] at (2.9,1.95) {$U$};
 \node at (-1.4,.64) {$10$}; \node at (1.4,.64) {$7$};
 \node at (0,1.17) {$5$}; \node at (-.76,-.77) {$4$};
 \node at (.76,-.77) {$3$}; \node at (0,-.17) {$3$};
 \node at (0,-1.64) {$6$}; \node at (2.52,-2.03) {$2$};
\end{tikzpicture}
\end{center}
Právě ve dvou: $5+4+3=12$.
Ve sjednocení: $10+7+6+5+4+3+3=38$;
mimo všechny $40-38=2$.
\textbf{Kontrola:} $10+7+6+5+4+3+3+2=40$ a
$|A\cap B|=5+3=8$.
\end{prikladReseniBox}

\begin{zadaniReseniBox}[title={R2 -- Intervaly kombinované se závorkami}]
Pro $I=(-\infty;3\rangle$, $J=\langle-1;5)$,
$K=(2;+\infty)$ urči $(I\cap J)\setminus K$ a $I\setminus(J\cup K)$.
\tcblower
$I\cap J=\langle-1;3\rangle$, proto
$(I\cap J)\setminus K=\langle-1;2\rangle$.
$J\cup K=\langle-1;+\infty)$, proto
$I\setminus(J\cup K)=(-\infty;-1)$.
\end{zadaniReseniBox}

\blok{Domácí procvičení před testem}
\begin{zadaniReseniBox}[title={D1 -- Zopakuj různé typy příkladů}]
\begin{enumerate}[label=\arabic*)]
\item Pro $U=\{0,1,\ldots,10\}$, $A=\{0,2,4,6,8,10\}$,
$B=\{3,6,9\}$ vypočti $A\cap B$, $A\setminus B$ a $(A\cup B)'_U$.
\item Urči průnik, sjednocení a oba rozdíly pro
$I=(-4;2\rangle$, $J=\langle2;6)$.
\item Rozhodni, zda $(-2;1)\cup(1;4)$ je interval.
Vysvětli, zda do něj patří $1$ a zda do něj patří $1{,}5$.
\item Ze $36$ studentů hraje šachy $20$, volejbal $17$ a ani jednu
z těchto aktivit nedělá $7$. Kolik se věnuje oběma aktivitám,
právě jedné a nejvýše jedné? Nakresli diagram.
\item Proč nelze bez určení základní množiny jednoznačně zapsat $A'$?
Uveď dva různé doplňky téže množiny $A=\{1,2\}$.
\end{enumerate}
\tcblower
\begin{enumerate}[label=\arabic*)]
\item $A\cap B=\{6\}$;
$A\setminus B=\{0,2,4,8,10\}$;
$(A\cup B)'_U=\{1,5,7\}$.
\item $I\cap J=\{2\}$;
$I\cup J=(-4;6)$;
$I\setminus J=(-4;2)$;
$J\setminus I=(2;6)$.
\item Není; číslo $1$ chybí, ačkoli $1{,}5$ do druhého intervalu patří.
\item Alespoň jedna aktivita: $36-7=29$. Obě: $20+17-29=8$.
Pouze šachy $12$, pouze volejbal $9$, právě jedna $21$,
nejvýše jedna $21+7=28$. Kontrola $12+8+9+7=36$.
\item Např. v $U_1=\{1,2\}$ je $A'_{U_1}=\varnothing$,
v $U_2=\{1,2,3\}$ je $A'_{U_2}=\{3\}$.
\end{enumerate}
\end{zadaniReseniBox}

\poznamka{Připomenutí pro další hodinu: Nejčastější chyby jsou
zaměnění $\in$ a $\subseteq$, opomenutí základní množiny u doplňku,
nesprávné závorky v krajních bodech a dvojí započtení průniku.
Příští hodinu následuje dvacetiminutový test z množin a poté úvod do sinu a kosinu.}
\end{document}
